\chapter{Conclusiones}
\label{chap:conclusiones}

Este capítulo cierra el trabajo recapitulando lo realizado, comprobando en qué medida
se han alcanzado los objetivos planteados y sintetizando las conclusiones que se
desprenden de la comparación experimental. A diferencia del capítulo de resultados,
que presenta y discute las cifras en detalle, aquí se adopta una lectura de conjunto:
qué enseña el estudio sobre la elección de un \textit{framework} de aprendizaje federado, qué
aporta respecto al estado del arte, con qué limitaciones debe interpretarse y qué
líneas quedan abiertas. No se introducen datos ni resultados nuevos; las conclusiones
se apoyan en la evidencia ya expuesta en los capítulos anteriores.

\section{Consecución de los objetivos}
\label{sec:consecucion-objetivos}

El objetivo general del trabajo era comparar de forma experimental varios \textit{frameworks}
de aprendizaje federado ejecutando un mismo escenario de entrenamiento en todos ellos
y observando tanto el rendimiento del modelo como el consumo de recursos y la
experiencia de uso. Este objetivo se ha cumplido: Flower, NVIDIA FLARE y FedML se han
implementado y evaluado sobre un escenario común —clasificación de CIFAR-10 con una
ResNet-18 adaptada, el algoritmo FedAvg y una configuración cross-silo horizontal— y
se han contrastado a lo largo de tres ejes complementarios. Los objetivos específicos
también se han satisfecho de la manera que se resume a continuación.

El primer objetivo específico, relativo al estudio de los fundamentos del aprendizaje
federado, se abordó en el capítulo de fundamentos teóricos, que delimitó el
funcionamiento de la ronda federada, el algoritmo FedAvg, los tipos de aprendizaje
federado y los retos derivados de la heterogeneidad de datos y recursos. El segundo,
centrado en la necesidad de \textit{frameworks} especializados y en los aspectos que influyen
en su elección, se cubrió en el estado del arte, donde se revisó el panorama de
herramientas, se fijaron los criterios de selección y se justificó la terna escogida.

El tercer y el cuarto objetivo, de diseño e implementación, se materializaron en el
capítulo de metodología y en el de diseño experimental e implementación: se definió un
escenario común con el conjunto de datos, el modelo, el algoritmo y los hiperparámetros
constantes, y se adaptó ese escenario a la arquitectura propia de cada herramienta. El
quinto objetivo,
relativo a la recogida de métricas de aprendizaje, rendimiento y consumo de recursos,
y el sexto, sobre la evaluación de la experiencia práctica de uso, se desarrollaron en
el capítulo de resultados mediante una comparativa cuantitativa de rendimiento y dos
valoraciones cualitativas, de usabilidad y de flexibilidad. Por último, el séptimo
objetivo, de comparación crítica y señalamiento de ventajas, limitaciones y escenarios
de uso, se atendió en la comparativa transversal entre \textit{frameworks} y se recoge,
condensado, en las conclusiones que siguen.

\section{Conclusiones principales}
\label{sec:conclusiones-principales}

La conclusión central del trabajo es que, bajo idénticas condiciones de datos, modelo
e hiperparámetros, la elección del \textit{framework} influye poco en la calidad del modelo y
mucho en el coste de obtenerlo, en la experiencia de desarrollo y en la estabilidad
del modo concurrente. Esta idea, que articula los tres ejes, se concreta a
continuación.

\subsection{Rendimiento}
\label{sec:concl-rendimiento}

En modo secuencial, los tres \textit{frameworks} alcanzaron precisiones muy próximas en el
experimento principal, en torno al 90,7--92,0\,\%, con una separación máxima de poco más
de un punto porcentual entre el mejor y el peor resultado. Esa convergencia confirma
que, cuando se mantienen constantes el conjunto de datos, la arquitectura del modelo y
la configuración de entrenamiento, la herramienta apenas condiciona la precisión final
del modelo global. Las diferencias relevantes aparecieron en otros indicadores: el
tiempo de ejecución, la utilización de la GPU y, sobre todo, el comportamiento en modo
concurrente. En el experimento principal, NVIDIA FLARE y FedML registraron tiempos
prácticamente equivalentes, ligeramente inferiores al de Flower, y FedML fue el único cuyo modo
concurrente llegó a completarse en el entorno utilizado, aunque esas ejecuciones no
proporcionaron una precisión final comparable en los registros analizados.
El modo concurrente marcó además el límite del equipo utilizado: ninguna de las tres
herramientas resultó inmune a las restricciones de recursos, si bien solo el backend
MPI de FedML representó una concurrencia real entre procesos.

\subsection{Usabilidad}
\label{sec:concl-usabilidad}

En el eje de usabilidad, los tres \textit{frameworks} resultaron funcionalmente aptos para el
caso de estudio, pero ofrecieron experiencias de desarrollo claramente distintas.
Flower fue la herramienta más accesible y de menor coste de aprendizaje, con una
puesta en marcha directa y un flujo secuencial estable. NVIDIA FLARE destacó por la
trazabilidad y la riqueza de sus registros, a costa de una instalación sensible a las
versiones, un número elevado de conceptos propios y una depuración exigente. FedML
ocupó una posición intermedia, equilibrada pero condicionada por la fiabilidad de su
instrumentación y por el trabajo adicional de extracción de métricas. De esta lectura
se desprende una recomendación dependiente del perfil del proyecto: Flower para
prototipado rápido, docencia o trabajos en los que el tiempo hasta el primer resultado
es crítico; NVIDIA FLARE cuando se priorizan la trazabilidad, la reproducibilidad
rigurosa y la cercanía a un despliegue real; y FedML cuando interesa disponer de varios
backends de ejecución bajo una misma configuración.

\subsection{Flexibilidad}
\label{sec:concl-flexibilidad}

En flexibilidad, las tres herramientas obtuvieron una valoración global equivalente, lo
que indica que todas permiten el tipo de adaptaciones que exige este trabajo: integrar
un modelo y un conjunto de datos propios, editar el entrenamiento local, sustituir la
estrategia de agregación, registrar métricas a medida y reproducir los experimentos
mediante ficheros y scripts. La diferenciación entre ellas no está, por tanto, en la
cifra global, sino en el perfil cualitativo: la simplicidad para sustituir la estrategia
de agregación y la orientación al despliegue en la nube de Flower; la profundidad en
seguridad y privacidad y la vocación de producción de NVIDIA FLARE; y la extensibilidad
por componentes y la experimentación algorítmica de FedML.

Un hallazgo transversal a los tres ejes es la fragilidad de las
estrategias adaptativas de la familia FedOpt. En los tres \textit{frameworks} fue posible
configurar FedAdam y FedAdagrad, pero su ejecución resultó inestable con poca
configuración —divergían, fallaban o no llegaban a aprender según el caso—, lo que
evidencia que la mera disponibilidad de una estrategia no equivale a su estabilidad:
esta depende de un ajuste deliberado de hiperparámetros, especialmente cuando el número
de rondas es reducido.

La lectura conjunta de los tres ejes conduce, así, a una conclusión única: seleccionar un
\textit{framework} de aprendizaje federado atendiendo solo a su rendimiento resulta insuficiente,
porque en condiciones equivalentes las tres herramientas rinden de forma prácticamente
indistinguible y la diferencia real reside en cuánto cuesta llegar a ese modelo, en lo
cómodo que resulta el desarrollo y en el margen de adaptación que ofrece cada una.

\section{Aportaciones del trabajo}
\label{sec:aportaciones}

La aportación principal de este trabajo es una comparación experimental, razonada y
reproducible de tres \textit{frameworks} de aprendizaje federado, que ayuda a cubrir un hueco
detectado en el estado del arte: la mayoría de los estudios previos se centra en
aspectos cualitativos o documentales y rara vez mide, sobre un mismo escenario y en
condiciones equivalentes, tanto la precisión del modelo como el consumo de recursos de
cada herramienta. Sobre esa idea general, el trabajo aporta varios elementos concretos.

En primer lugar, el diseño de un escenario experimental común y controlado, con el
conjunto de datos, el modelo, el algoritmo, los hiperparámetros y el hardware fijados,
que permite atribuir las diferencias observadas al \textit{framework} y no a otros factores. En
segundo lugar, una evaluación articulada en tres ejes complementarios —rendimiento,
usabilidad y flexibilidad— que ofrece una visión más completa que la habitual centrada
solo en la precisión o en la velocidad. En tercer lugar, dos instrumentos de valoración
reutilizables: una rúbrica de usabilidad basada en las dimensiones de experiencia de
uso y una rúbrica de flexibilidad con diez criterios, que separan de forma explícita la
evidencia experimental de la documental. Por último, la documentación rigurosa del
comportamiento en modo concurrente, incluidas las incidencias de cada herramienta, que
constituye una evidencia práctica de interés para quien deba afrontar este tipo de
ejecuciones en hardware limitado.

\section{Limitaciones del estudio}
\label{sec:limitaciones-conclusiones}

Las conclusiones anteriores deben interpretarse a la luz de las limitaciones ya
señaladas en la metodología. La principal es de hardware: la GPU de 6\,GB y los 14\,GiB
de memoria RAM del equipo condicionaron el modo concurrente en las tres herramientas, de
modo que el comportamiento concurrente refleja tanto el diseño de cada \textit{framework} como el
límite del entorno de pruebas. Además, cada herramienta se ejecutó en un entorno de
software distinto, con versiones diferentes de Python, PyTorch y CUDA, lo que supone una
amenaza menor a la validez de la comparación de rendimiento, mitigada al documentar las
versiones exactas.

A ello se añaden otras restricciones de alcance. El estudio se limita a una distribución
de datos IID, por lo que no evalúa el efecto de la heterogeneidad no IID. Los
hiperparámetros se fijaron a priori, sin optimización individual por \textit{framework}, de modo
que los resultados corresponden a esa configuración concreta. Cada configuración se
ejecutó una sola vez con una semilla fija, por lo que no se dispone de medidas de
variabilidad entre repeticiones y los valores deben leerse como una única observación
por configuración. Toda la experimentación se realizó en simulación sobre una sola
máquina, sin latencias de red reales ni efectos propios de un despliegue distribuido.
Por último, las métricas de consumo de recursos dependen del hardware empleado y las
valoraciones de usabilidad y flexibilidad, aunque apoyadas en evidencias registradas,
incorporan el juicio del autor como desarrollador, por lo que deben tomarse como una
guía cualitativa y no como una medida absoluta.

\section{Líneas de trabajo futuro}
\label{sec:trabajo-futuro}

El trabajo abre varias líneas de continuación que prolongan de forma natural lo
realizado.

\begin{itemize}
  \item \textbf{Escenarios no IID.} Repetir la comparación con particiones no
  independientes ni idénticamente distribuidas permitiría observar cómo afecta la
  heterogeneidad de los datos al rendimiento y a la convergencia de cada \textit{framework}, un
  factor deliberadamente aislado en este estudio.

  \item \textbf{Recuperación de la precisión del modo concurrente de FedML.} Dado que
  las ejecuciones concurrentes de FedML completaron pero no registraron una precisión
  final extraíble, una evaluación a posteriori del modelo guardado, con el mismo
  procedimiento de evaluación centralizada del modo secuencial, permitiría incorporar
  esos resultados a la comparación.

  \item \textbf{Análisis de variabilidad.} Ejecutar cada configuración con varias
  semillas haría posible reportar media y desviación típica y estimar la significación
  de las diferencias observadas, superando la limitación de la observación única.

  \item \textbf{Ajuste de las estrategias adaptativas.} Profundizar en la
  configuración de la familia FedOpt —tasa de aprendizaje del servidor, factor de
  adaptabilidad y número de rondas— permitiría determinar en qué condiciones FedAdam y
  FedAdagrad resultan competitivos frente a FedAvg.

  \item \textbf{Despliegue distribuido real.} Trasladar el experimento de la simulación
  a un despliegue con clientes en máquinas distintas mediría latencias de red y efectos
  de un entorno federado realista, especialmente relevante para NVIDIA FLARE y FedML,
  orientados a producción.

  \item \textbf{Mecanismos de seguridad y privacidad.} Evaluar de forma experimental la
  privacidad diferencial, la agregación segura o el cifrado, capacidades que en este
  trabajo solo se verificaron documentalmente, completaría el análisis de flexibilidad
  con evidencia ejecutada.
\end{itemize}

\section{Reflexión final}
\label{sec:reflexion-final}

El desarrollo de este trabajo ha mostrado que comparar herramientas de aprendizaje
federado exige mucho más que medir precisiones: obliga a implementar un mismo escenario
sobre arquitecturas muy distintas, a instrumentar la ejecución con cuidado y a
distinguir con honestidad entre lo que se ha verificado experimentalmente y lo que solo
consta en la documentación. El resultado es una fotografía equilibrada de tres
herramientas que, lejos de poder ordenarse en un ranking absoluto, responden a
prioridades complementarias. La principal enseñanza, tanto para este estudio como para
quien deba seleccionar un \textit{framework} en un proyecto similar, es que la decisión debe
guiarse por el contexto —el coste asumible, la experiencia de desarrollo buscada y las
necesidades de flexibilidad y despliegue— y no por una supuesta superioridad en
precisión que, en condiciones equivalentes, apenas distingue a unas herramientas de
otras.
